\begin{cor}[Corollary B: the conifold gap of the formal quintic]
\label{cor:mai-gap}
Let $g\geq2$. In every conifold frame of \cref{def:set-frames},
\[
F_g^{\mathrm{cf}}\in c_g\,\tc^{2-2g}+\bbQ\fps{u}.
\]
Since $\tc=u+O(u^2)$ lies in $\bbQ\fps{u}$, $\bbQ\fps{u}=\bbQ\fps{\tc}$: the conifold expansion of $F_g$ contains no fractional power of $u$ and no power $\tc^k$ with $3-2g\leq k\leq-1$, and its leading term is $c_g\tc^{2-2g}$.
\end{cor}

\begin{proof}
Write $\Cexp$ for $\Cexp_{\alpha,\beta}$, extended $\bbQ(\zeta)$-linearly.

\emph{Step 1: expanding the graph sum.} By \cref{lem:gau-confexp}, $\Cexp$ is a ring homomorphism, and $F_g^{\Nring}$ is a polynomial expression in finitely many elements of $\Nring\otimes\bbQ(\zeta)$: the coefficients of the edges of \cref{def:mai-Fcf}, of the leaf points and of $\TFT$; and the edge $E_X$, a polynomial expression in the coefficients of $X^{-1}$ followed by an exact division by $\psi+\psi'$, commutes with ring homomorphisms. Hence
\[
F_g^{\mathrm{cf}}=\sum_\Gamma\mathrm{Cont}_\Gamma\bigl[E_{\Cexp(\tilde R)},q_{\Cexp(\tilde R)};\Cexp(\TFT)\bigr].
\]
Since $\Dn=1$ (\cref{thm:gam-D}), \cref{thm:gau-R0}~(iv) and \cref{cor:gau-colours} give the identity $\Cexp(\Qr)=\Pc\Qr^{\mathrm m}$ for the expanded rows. Here the canonical gauge is symplectic, $\Pc(-z)^T\eta\,\Pc(z)=\eta$ (\cref{prop:gau-symp}), and so $\Pc(-z)^T\eta^{-1}\Pc(z)=\eta^{-1}$, as $\eta^{-1}=\eta/25$; and the model columns satisfy $\Qr^{\mathrm m}(-z)^T\eta^{-1}\Qr^{\mathrm m}(z)=1$, because $\Qr^{\mathrm m}=\Lambda_\varsigma^{-1}\Psi_0\Mo(\tau\zh)$ (\cref{lem:flh-inputs}~(i)) and
\[
\Lambda_\varsigma^{-1}\eta^{-1}\Lambda_\varsigma^{-1}=\eta^{-1},\qquad \Psi_0^T\eta^{-1}\Psi_0=1,\qquad \Mo(-t)^T\Mo(t)=1,
\]
the first because $\nu_a+\nu_{3-a}=2$, the second by \cref{def:set-M0}, the third by \cref{lem:set-M0sympl}. Hence $\Cexp(\Qr)(-z)^T\eta^{-1}\Cexp(\Qr)(z)=1$. With $\Psi_{\mathrm{cf}}=\Cexp(\Psi_{\mathrm Q})$, so that $\Psi_{\mathrm{cf}}^T\eta^{-1}\Psi_{\mathrm{cf}}=1$ and $\Psi_{\mathrm{cf}}\Psi_{\mathrm{cf}}^T=\eta$ (\cref{lem:mai-Psi}), and $\Cexp(\tilde R)=\Psi_{\mathrm{cf}}^{-1}\Cexp(\Qr)$, this gives
\[
\Cexp(\tilde R)(-z)^T\,\Cexp(\tilde R)(z)=\Cexp(\Qr)(-z)^T(\Psi_{\mathrm{cf}}\Psi_{\mathrm{cf}}^T)^{-1}\Cexp(\Qr)(z)=1 :
\]
$\Cexp(\tilde R)$ is symplectic, and the edge $E_{\Cexp(\tilde R)}$ of \cref{def:mai-Fcf} is the edge of \cite[Def.~2.2]{PPZ15}. The classes of $\Cexp(\TFT)$ are those of $\TFT$ with $\Delta_i^{\pm1/2}$ replaced by $\Cexp(\Delta_i^{\pm1/2})$, and the axioms of a topological field theory with unit are polynomial identities in the $\Delta_i^{\pm1/2}$; so $\Cexp(\TFT)$ is a topological field theory over $\bbQ(\zeta)\lser{s}$ with the unit $\Cexp(1_{\TFT})$, and $q_{\Cexp(\tilde R)}=\Cexp(\tilde R)^{-1}(z)\bigl(-z\,\Cexp(1_{\TFT})\bigr)$. By \cref{lem:set-pushins}~(iii), the right-hand side above is therefore $\int_{\overline{\mathcal M}_g}\bigl(\Cexp(\tilde R).\Cexp(\TFT)\bigr)_{g,0}$, the integral of the genus-$g$ class without legs of the unit-preserving action of \cite[Def.~2.13]{PPZ15}.

\emph{Step 2: regrouping.} As a series in $z$ over $\bbQ(\zeta)\lser{s}$, $M=\Mo(\tau(u)\zh)$ satisfies $M(0)=1$ and is symplectic. Then $X\coloneqq\Cexp(\tilde R)M^{-1}$ is symplectic with $X(0)=1$, and $\Cexp(\Qr)M^{-1}=\Psi_{\mathrm{cf}}X$ is the outer matrix $R'$ of \cref{def:set-outer}. By \cite[Prop.~2.14]{PPZ15}, the unit-preserving action of \cite[Def.~2.13]{PPZ15} is a left group action, so
\[
\Cexp(\tilde R).\Cexp(\TFT)=X.\bigl(M.\Cexp(\TFT)\bigr).
\]
(The results of \cite{PPZ15} are stated over $\bbQ$; the proofs of their Props.~2.4, 2.12 and 2.14, and of Props.~2.8 and~2.9, which the proof of Prop.~2.14 uses, are identities between finite sums, polynomial in the coefficients of the $R$-matrices and of the CohFT, and hold verbatim over $\bbQ(\zeta)\lser{s}$.) Take the genus-$g$ parts without legs and integrate over $\overline{\mathcal M}_g$. The left-hand side gives $F_g^{\mathrm{cf}}$ by Step 1. The right-hand side gives $\sum_\Gamma\mathrm{Cont}_\Gamma[E_X,q_X;M.\Cexp(\TFT)]$ by \cref{lem:set-pushins}~(iii), applied to the cohomological field theory $M.\Cexp(\TFT)$, whose unit is that of $\Cexp(\TFT)$ by \cite[Prop.~2.12]{PPZ15}, and to the symplectic series $X$ with $X(0)=1$; this is the contraction of \cref{def:flh-outer} with $\varkappa=1$. Since $\Psi_{\mathrm{cf}}$ is an isometry from the colour space onto $(\bbH,\eta)$ independent of $z$ which maps $1_{\TFT}$ to $1$ (\cref{lem:mai-Psi}~(ii), \cref{lem:flh-inputs}~(iv)), $R'(-\psi)^T\eta^{-1}R'(-\psi')=X(-\psi)^TX(-\psi')$ and $R'^{-1}(z)(-z1)=X^{-1}(z)(-z1_{\TFT})$; so $E_{R'}=E_X$ and $q_{R'}=q_X$, by the computation of the proof of \cref{lem:flh-frame} for the isometry $A=\Psi_{\mathrm{cf}}$. Each term is therefore $\mathrm{Cont}_\Gamma[E_{R'},q_{R'};M.\TFT]=T_\Gamma$, with $\TFT$ expanded in the frame as in \cref{def:set-outer}, and
\[
F_g^{\mathrm{cf}}=\sum_\Gamma T_\Gamma ,
\]
with no further factor $|\Aut\Gamma|^{-1}$.

\emph{Steps 3 and 4.} By Theorem~A (\cref{thm:mai-main}), $T_\Gamma\in\bbQ\fps{u}$ for every $\Gamma$ with an edge, and the vertex graph gives $c_g\tc^{2-2g}+\bbQ\fps{u}$. The sum is finite, so $F_g^{\mathrm{cf}}\in c_g\tc^{2-2g}+\bbQ\fps{u}$.
\end{proof}

Corollary~B uses Theorem~A, the reconstruction (T) through the identification $\ev(F_g^{\Nring})=F_g$, \cite[\S0.2, (1); Def.~0.3; (5); Rem.~2.1; Defs.~2.2 and~2.5, with the conventions stated after them; (8); Prop.~2.12; Def.~2.13; Prop.~2.14, with Props.~2.4, 2.8 and~2.9 through its proof]{PPZ15}, and proved statements only (\cref{lem:set-pushins}, \cref{lem:gau-confexp}, \cref{cor:gau-rowsNN}, \cref{lem:gau-colrows}, \cref{thm:gau-R0}, \cref{cor:gau-colours}, \cref{thm:gam-D}, \cref{prop:gau-symp}, \cref{lem:set-M0sympl}, \cref{lem:mai-Psi}, \cref{lem:flh-inputs}, \cref{lem:flh-frame}); its conditionality is that of Theorem~A.

\begin{rmk}[the constant]
\label{rmk:mai-constant}
The constant $c_g$ is the constant of \cite[Conj.~1.5]{CGYZ26}, which is proved there for the quintic threefold \cite[Thm.~2]{CGYZ26}, in the same normalisation of the vanishing period; and the two conifold expansions agree on the periods common to both settings. The values of $c_2,\dots,c_5$ are in \cref{rmk:set-cg}.

\emph{The dictionary.} In the notation of the note, \cite[Def.~1.1]{CGYZ26} works over $\bbQ(q)$ with $\Dq=q\,\dd/\dd q$, the periods $C_0$ and $I_1$ of $\PFfour$ (the formal quintic has the genus-$0$ periods of the quintic, \cref{def:set-genus0}), the mirror map $T=I_1/C_0$, and the differential $\bbQ(q)$-algebra generated by $C_0^{\pm1}$ and $C_1^{\pm1}$, $C_1=\Dq T$. A pair of periods at the conifold point, one not vanishing there and one vanishing simply, with a logarithmic period equal to the vanishing one times $\log\Delta$ plus a regular series, is admissible \cite[Def.~1.2]{CGYZ26} if the substitution of these periods for $C_0$ and of the derivative of their quotient for $C_1$ respects the relations of that algebra; it defines a differential homomorphism to Laurent series in $\Delta=1-5^5q$, the conifold expansion of \cite{CGYZ26}. Their example \cite[Ex.~1.3]{CGYZ26} is the pair $(f_0,\varpi_1)$ of \cref{lem:set-frobbasis}, i.e. the frame $(\alpha,\beta)=(0,0)$ with $\Delta=u$ and their flat coordinate equal to $\tc$; their family \cite[Rem.~1.4]{CGYZ26} is exactly \cref{def:set-frames}, with $(\alpha,\beta)\in\bbQ^2$; and the normalisation $\varpi_1=u+O(u^2)$ is that of \cite[\S4.2, Prop.-Def.~4.5\,(1)]{CGYZ26}. Their expansion and $\Cexp_{\alpha,\beta}$ are differential homomorphisms for the same $\Dq$ which both send $C_0$ to $\varpi_0$ and $C_1$ to $\Dq\tc$ and both expand at $u=0$, where $q=(1-u)/5^5$. So they agree on the differential subring of $\Nring$ generated by $c_0^{\pm1}$ and $(\mathsf W/c_0^2)^{\pm1}$ over $\bbQ[L^5,L^{-5}]=\bbQ[q,(1-5^5q)^{-1}]$, wherever theirs is defined. The constant of \cite[Conj.~1.5]{CGYZ26} is $(-1)^{g-1}B_{2g}/(5^{g-1}\,2g(2g-2))=c_g$. The comparison is with that conjecture only; the sign convention of \cite[Def.~1.8]{CGYZ26} is different and is not used.

\emph{The normalisation.} Only the scale of $\varpi_1$ matters. Let $\varkappa_0,\varkappa_1\in\bbC^\times$ (local letters), let $\Cexp'$ be the conifold expansion of \cref{def:set-confexp} with $(\varpi_0,\varpi_1)$ replaced by $(\varkappa_0\varpi_0,\varkappa_1\varpi_1)$, and $\tc'\coloneqq(\varkappa_1/\varkappa_0)\tc$. Then
\[
\Cexp'(F_g^{\Nring})-c_g\,\varkappa_1^{2g-2}\,\tc'^{\,2-2g}\in\bbC\fps{u}:
\]
the constant does not depend on $\varkappa_0$, and it is $c_g$ if and only if $\varkappa_1^{2g-2}=1$. Indeed, the $\bbC$-linear endomorphism $\mathfrak k$ (a local letter) of $\Nring\otimes\bbC$ with $c_j\mapsto\varkappa_0c_j$, $i_j\mapsto\varkappa_1i_j+(\varkappa_1-\varkappa_0)(\logq\,c_j+jc_{j-1})$, $\logq\mapsto\logq$, $L\mapsto L$ commutes with $\Dq$ ($\PFfour$ is linear and the concomitant relation bilinear) and satisfies $\Cexp'=\Cexp\circ\mathfrak k$; it sends $C_0$, $I_1$, $C_1$, $\mathsf W$ to $\varkappa_0C_0$, $\varkappa_1I_1$, $(\varkappa_1/\varkappa_0)C_1$, $\varkappa_0\varkappa_1\mathsf W$. It fixes the flat amplitudes $A_n$, which lie in $\bbQ(\zeta)[L,L^{-1}]$ (\cref{thm:gau-R0}~(iii); in $\bbQ[L,L^{-1}]$ by \cref{thm:con-glue}), so by the forward relations of \cref{prop:set-cyclicLR}, whose coefficients $C_1\kappa^{\mathrm{LR}}_a$ are $(C_1,L^5/(C_0^2C_1^2),C_1,C_0,C_0)$, it multiplies the rows by $\diag(\varkappa_0^{-1},\varkappa_1^{-1},\varkappa_1,\varkappa_0,1)$, which preserves $\eta^{-1}$. It multiplies $\Psi_{\mathrm Q}$ by the same matrix, so $\mathfrak k(\tilde R)=\tilde R$, while $\Delta_i^{1/2}\mapsto\varkappa_0\Delta_i^{1/2}$ and $1_{\TFT}\mapsto\varkappa_0^{-1}1_{\TFT}$. Hence $\mathfrak k$ fixes the edges, multiplies every leaf insertion by $\varkappa_0^{-1}$ and every vertex with $n_v$ half-edges and $l_v$ leaf insertions by $\varkappa_0^{2g_v-2+n_v+l_v}$; as $\sum_v(2g_v-2+n_v+l_v)=2g-2+\sum_vl_v$, every term is multiplied by $\varkappa_0^{2g-2}$, which is the count in the proof of \cref{lem:flh-comm}. So $\mathfrak k(F_g^{\Nring})=\varkappa_0^{2g-2}F_g^{\Nring}$, and by Corollary~B, $\Cexp'(F_g^{\Nring})=\varkappa_0^{2g-2}F_g^{\mathrm{cf}}=c_g\varkappa_1^{2g-2}\tc'^{\,2-2g}+\varkappa_0^{2g-2}\bbQ\fps{u}$. Consequently, for every pair of the family of \cite[Rem.~1.4]{CGYZ26} with the vanishing period normalised as in \cite[Prop.-Def.~4.5\,(1)]{CGYZ26}, whatever the scale of the other period, the formal quintic has the conifold gap with the constant of \cite[Conj.~1.5]{CGYZ26}. The conventions of the literature for this constant are collected in \cref{sec:sgn-lit}.

\emph{What is not claimed.} (1) The two theories are different: the potentials of the formal quintic are not those of the quintic, and they differ already in the constant term at genus $2$ (\cref{sec:evi-calib}). The equality of the constants concerns two different potentials; nothing about the quintic follows. (2) The rows of the formal quintic involve $L$, which is not in $\bbQ(q)$: $\Nring$ contains $L^{\pm1}$, and $\Cexp(L)=s^{-1}$ is not a Laurent series in $u$. The total $F_g^{\mathrm{cf}}$ is nevertheless one, by Corollary~B. Whether $F_g^{\Nring}$ lies in the algebra of \cite[Def.~1.1]{CGYZ26} is neither needed nor proved. (3) The expansion of \cite{CGYZ26} is defined on functions, under the admissibility assumption; $\Cexp$ is defined on the presented ring $\Nring$, with no assumption (\cref{lem:gau-confexp}), and the conifold expansion of $F_g$ is that of its representative $F_g^{\Nring}$ (\cref{def:mai-Fcf}). (4) The degree-$0$ term of $F_g$ is a constant and does not affect the statement. (5) \cite[Rem.~1.9]{CGYZ26} names the conifold gap of the formal quintic as one of the two key statements of the approach of Guo, Janda and Ruan \cite{GJR18}, without stating it; Corollary~B is not claimed to be literally that statement.
\draftnote{Numbers from arXiv:2609.13921v2, read in the HTML rendering; to be re-read in the PDF.}
\end{rmk}

\subsection{Consequences of \texorpdfstring{$\Dn=1$}{D=1} for the natural form}
\label{sec:mai-nat}

By \cref{thm:gam-D} the normalisation series is trivial. This subsection records what this means for the natural form: the hatted rows factor through the model's $R$-matrix in its natural variable (\cref{cor:mai-natfact}), and the actual gauge is the canonical one (\cref{cor:mai-actual}). None of it is used in the proof of Theorem~A.

Let $\hat\Qr(u,\zh)=\Lambda_s\Qr(u,s^4\zh)$ be the hatted rows of \cref{def:set-hat} (flat index $a$, colour $i$). Its constant term is $\hat\Qr(u,\zh=0)=\Lambda_s\Psi_{\mathrm{cf}}=\Psi(u)$, the matrix of \cref{lem:mai-Psi}~(ii), and we put
\[
\tilde R(\zh;u)\coloneqq\Psi(u)^{-1}\hat\Qr(u,\zh)=\Psi(u)^T\eta^{-1}\hat\Qr(u,\zh).
\]
Since $\Psi(u)^{-1}\hat\Qr(u,\zh)=\Psi_{\mathrm{cf}}^{-1}\Cexp(\Qr)(s^4\zh)$, $\tilde R(\zh;u)$ is the conifold expansion of the matrix $\tilde R$ of \cref{def:mai-Fcf}, written in the variable $\zh$.

\begin{lem}[the hatted gauge]
\label{lem:mai-Ghat}
Write the graded canonical gauge as $\Pc=\sum_kz^k\sum_bp_{k,b}(u)E_{b+k,b}$ with $p_{k,b}\in\bbQ\fps{u}$ (\cref{thm:gau-U}, \cref{prop:gau-graded}, \cref{def:gau-can}), and put $\Ghat(u,\zh)\coloneqq\Lambda_s\Pc(u,s^4\zh)\Lambda_\varsigma^{-1}$. Then
\[
\Ghat_{ab}=\tau^{(1-\nu_b)/4}\sum_k\zh^k\,u^{\mu(b,k)}\,p_{k,b}(u)\,\delta_{a,b+k},\qquad\mu(b,k)\coloneqq\frac{4k+\nu_b-\nu_{b+k}}{5},
\]
where $\mu(b,k)$ is an integer with $\mu(b,k)\geq\max(0,(4k-4)/5)$, and $\mu(b,k)=0$ only for $k=0$ and for $(b,k)=(1,1)$, the entry $(2,1)$. Hence $\Ghat$ is a power series in $u$ whose coefficient of $u^n$ is a polynomial in $\zh$ of degree at most $\lfloor(5n+4)/4\rfloor$.
\end{lem}

\begin{proof}
Since $\varsigma=s\tau^{-1/4}$ (\cref{def:set-tau}), the entry $(b+k,b)$ of $\Lambda_s\Pc(u,s^4\zh)\Lambda_\varsigma^{-1}$ is $p_{k,b}\zh^k\tau^{(1-\nu_b)/4}$ times $s$ to the power $(1-\nu_{b+k})+4k-(1-\nu_b)=4k+\nu_b-\nu_{b+k}$. As $\nu_a\equiv-a\bmod5$, this exponent is $\equiv5k\equiv0\bmod5$, so it is $5\mu(b,k)$ and $s^{5\mu(b,k)}=u^{\mu(b,k)}$. The $\nu_a$ take values in $\{-1,\dots,3\}$, so $\nu_b-\nu_{b+k}\geq-4$ and $\mu(b,k)\geq(4k-4)/5$; for $k=0$ the exponent is $0$, and for $k\geq1$ it is $\geq0$. For $k\geq2$ it is positive; for $k=1$ it vanishes iff $\nu_{b+1}-\nu_b=4$, which happens only for $b=1$. At a given $u^n$ only the $k$ with $(4k-4)/5\leq n$ occur.
\end{proof}

\begin{cor}[the natural form after $\Dn=1$]
\label{cor:mai-natfact}
In every conifold frame of \cref{def:set-frames}:
\begin{enumerate}
\item[(a)] $\hat\Qr(u,\zh)=\Ghat(u,\zh)\,\Rhat^{\mathrm m}(\tau(u)\zh)$. In particular $\hat\Qr$ is a power series in $u$ whose coefficients are power series in $\zh$, and $\Ghat$ is a power series in $u$ whose coefficients are polynomials in $\zh$ (\cref{lem:mai-Ghat}).
\item[(b)] $\hat\Qr(0,\zh)=\Rhat^{\mathrm m}(\zh)$: the conifold symbol of the formal quintic is that of the model. Hence $\tilde R(\zh;0)=\Mo(\zh)$ and $\Psi(0)=\Psi_0$.
\item[(c)] $\tilde R(\zh;u)=G(u,\zh)\,\Mo(\tau(u)\zh)$, with $\Mo$ on the right, where
\[
G\coloneqq\Psi(u)^{-1}\,\Ghat\,\Psi_0=\Psi_0^{-1}\,\Ghat(u,0)^{-1}\Ghat(u,\zh)\,\Psi_0 .
\]
Moreover $G(u,0)=G(0,\zh)=1$, the coefficient of every $u^n$ in $G$ is a polynomial in $\zh$, and $G(-\zh)^TG(\zh)=1$. The factor $\Mo(\tau\zh)$ is the same in every conifold frame; $G$ depends on the frame.
\end{enumerate}
\end{cor}

\begin{proof}
(a) By \cref{thm:gau-R0}~(iv) with \cref{cor:gau-colours} and $\Dn=1$, $\Qr=\Pc\Qr^{\mathrm m}$ for the expanded rows, and $\Qr^{\mathrm m}=\Lambda_\varsigma^{-1}\Rhat^{\mathrm m}(z/\varsigma^4)$ with $z/\varsigma^4=\tau\zh$ (\cref{prop:set-modrows}, \cref{def:set-tau}, \cref{lem:set-canon}; \cref{lem:flh-inputs}~(i)). Multiplying by $\Lambda_s$ and inserting $\Lambda_\varsigma^{-1}\Lambda_\varsigma$ gives $\hat\Qr=\Ghat\,\Rhat^{\mathrm m}(\tau\zh)$. The coefficients of $\Rhat^{\mathrm m}(\tau\zh)$ lie in $\bbQ(\zeta)\fps{u}$, so the polynomiality of $\Ghat$ (\cref{lem:mai-Ghat}) makes $\hat\Qr$ a power series in $u$ with coefficients in $\bbQ(\zeta)\fps{\zh}$; in particular $\hat\Qr$ can be evaluated at $u=0$.

(b) At $\zh=0$, (a) gives $\Psi(u)=\Ghat(u,0)\Psi_0$, since $\Rhat^{\mathrm m}(0)=\Psi_0$ ($\Fs_b(0)=\varepsilon_b$, \cref{prop:set-modrows}); with \cref{lem:mai-Psi}, $\Ghat(0,0)=\Psi(0)\Psi_0^{-1}=1$. By \cref{lem:mai-Ghat} the only entry of $\Ghat$ at $u^0$ off the diagonal is the entry $(2,1)$, with coefficient $\zh\,p_{1,1}(0)$, and $p_{1,1}(0)=0$ by \cref{lem:gau-R1}: column $1$ of $\Pc(0,z)$ is $(\Pc)_{11}(0,z)e_1$. So $\Ghat(0,\zh)=1$, and $\tau(0)=1$ gives $\hat\Qr(0,\zh)=\Rhat^{\mathrm m}(\zh)$. Then $\tilde R(\zh;0)=\Psi_0^{-1}\Rhat^{\mathrm m}(\zh)=\Mo(\zh)$ by \cref{def:set-M0}.

(c) $\tilde R=\Psi(u)^{-1}\Ghat\,\Rhat^{\mathrm m}(\tau\zh)=[\Psi(u)^{-1}\Ghat\,\Psi_0][\Psi_0^{-1}\Rhat^{\mathrm m}(\tau\zh)]=G\,\Mo(\tau\zh)$. The second expression for $G$ follows from $\Psi(u)=\Ghat(u,0)\Psi_0$; $\Ghat(u,0)=\diag(\tau^{(1-\nu_b)/4}p_{0,b})$ is invertible over $\bbQ\fps{u}$, the $p_{0,b}$ being units (\cref{thm:gau-U}), so every $u^n$-coefficient of $G$ is a polynomial in $\zh$. $G(u,0)=1$ is clear, and $G(0,\zh)=1$ is (b). For the symplecticity, $\Qr=\Pc\Qr^{\mathrm m}$ with $\Pc(-z)^T\eta^{-1}\Pc(z)=\eta^{-1}$ (\cref{prop:gau-symp}; $\eta^{-1}=\eta/25$) and $\Qr^{\mathrm m}(-z)^T\eta^{-1}\Qr^{\mathrm m}(z)=1$ (\cref{lem:set-M0sympl}, $\Psi_0^T\eta^{-1}\Psi_0=1$, and $\Lambda_\varsigma^{-1}\eta^{-1}\Lambda_\varsigma^{-1}=\eta^{-1}$) give $\Qr(-z)^T\eta^{-1}\Qr(z)=1$ (the case $\Dn=1$ of \cref{lem:flh-inputs}~(iii)), hence $\hat\Qr(-\zh)^T\eta^{-1}\hat\Qr(\zh)=1$ because $\Lambda_s\eta^{-1}\Lambda_s=\eta^{-1}$, and $\tilde R(-\zh)^T\tilde R(\zh)=1$ because $\Psi(u)^{-T}\Psi(u)^{-1}=\eta^{-1}$ (\cref{lem:mai-Psi}). Therefore $G(-\zh)^TG(\zh)=\Mo(-\tau\zh)^{-T}\Mo(\tau\zh)^{-1}=1$ by \cref{lem:set-M0sympl}. Neither $\Mo$ nor $\tau$ involves the frame.
\end{proof}

None of these conclusions is a test of the full identity $\Dn=1$. If $\Dn$ were a scalar series $\cn(z)$ in $1+z^5\bbQ\fps{z^5}$, the proof would give (a) and (b) with $\Pc$ replaced by the graded regular gauge $\Pc\cn(z)^{-1}$, and the polynomiality of the factor $\tilde R\,\Mo(\tau\zh)^{-1}$, which would then be $G\,\cn(z)^{-1}$, because $\cn(s^4\zh)^{-1}$ is a polynomial in $\zh$ at every order in $u$ ($z^5=u^4\zh^5$); and this factor is symplectic if and only if $\cn(z)\cn(-z)=1$. So only the symplecticity of the factor sees $\Dn$, and it sees only $\Dn_0(z)\Dn_0(-z)$, that is the even part of $\log\Dn_0$; a factor $\exp(a(z))$ with $a$ odd is invisible to all of (a)--(c) (compare \cref{rmk:con-why,rmk:mai-stored}, \cref{lem:flh-inputs}~(iii)). In the computations of \cref{sec:evi-normal} with rows built from a shifted free coefficient, which are not symplectic, the symplecticity of the factor is the first conclusion that fails; the same computations show that $\Mo$ cannot be placed on the left of $G$. Neither statement is used.

\begin{cor}[the actual gauge is canonical]
\label{cor:mai-actual}
Let $P^{\mathrm{act}}\coloneqq\Qr\,(\Qr^{\mathrm m})^{-1}$, the matrix series that carries the model columns at $v(u)$ to the rows of the formal quintic, both with their exponentials stripped (\cref{thm:gau-R0}). Then $P^{\mathrm{act}}=\Pc$. In particular $P^{\mathrm{act}}$ is a graded regular gauge with coefficients in $\bbQ\fps{u}$, $P^{\mathrm{act}}_{11}(0,z)=1$, and $P^{\mathrm{act}}\YC(v(u))=\YCQ$.
\end{cor}

\begin{proof}
$\Qr^{\mathrm m}$ is invertible because $\Qr^{\mathrm m}(0)=\Psi^{\mathrm m}$ is (\cref{lem:flh-inputs}~(i)). By \cref{thm:gau-R0}~(iv) with \cref{cor:gau-colours}, $\Qr\Dn=\Pc\Qr^{\mathrm m}$, and $\Dn=1$ by \cref{thm:gam-D}; so $P^{\mathrm{act}}=\Pc$. The remaining assertions are \cref{def:gau-can} and \cref{thm:gau-G}.
\end{proof}

With \cref{thm:gam-D} this is Theorem~C (\cref{thm:int-D}): the $R$-matrix of the formal quintic, normalised at $q=0$ by (QRR), is in the canonical normalisation at the conifold point. The two normalisations are imposed at different points, the rows at $u=1$ and the canonical gauge at $u=0$; their agreement is the content of \cref{thm:gam-D}, so the statement is not circular.

\begin{rmk}[a priori regular graphs]
\label{rmk:mai-apriori}
By \cref{cor:mai-natfact}~(c), the hatted outer matrix $\hat R'=\hat\Qr\,\Mo(\tau\zh)^{-1}$ equals $\Psi(u)\,G(u,\zh)$, with $G(0,\zh)=1$ and $\Psi(u)$ an isometry independent of $\zh$ that maps the unit of $\hat\TFT$ to $\hat1$ (\cref{lem:mai-Psi}, \cref{lem:flh-inputs}~(iv)). Hence the hatted edges and leaf points are those of $G(u,\zh)$: the edges are $(1-G(-\psi)^TG(-\psi'))/(\psi+\psi')$ and the leaf insertions are $\zh(1_{\hat\TFT}-G^{-1}(\zh)1_{\hat\TFT})$, $1_{\hat\TFT}$ the unit of $\hat\TFT$; both are $O(u)$. The hatted vertex CohFT $\Mo(\tau\zh).\hat\TFT$ has coefficients in $\bbQ(\zeta)\fps{u}$, its genus-$0$ three-point values $\zeta^{4i}\varpi_0$ being units there. So for a graph with $E$ edges every term of $G_\Gamma$ with $m$ leaf insertions is a power series in $u$ of order at least $E+m$, and by \cref{lem:flh-hat}, $T_\Gamma=u^{2-2g}G_\Gamma$ is $O(u^{E+2-2g})$. In particular every graph with $E\geq 2g-2$ is regular without further argument, and for the others only the coefficients $\coef{u^k}G(u,\zh)$, $1\leq k\leq 2g-3$, enter the singular part of $T_\Gamma$, besides the vertex classes modulo $u^{2g-2}$. This is not a step of the proof of Theorem~A.
\end{rmk}

\begin{rmk}[the rescaling $\tau$]
\label{rmk:mai-tau}
One expects $\tau$ to be the only rescaling for which the factor of \cref{cor:mai-natfact}~(c) is polynomial at every order in $u$: if, for some $\vartheta\in\bbQ\fps{u}$ with $\vartheta(0)=1$ (a local letter), every $u^n$-coefficient of $\tilde R\,\Mo(\vartheta\zh)^{-1}$ is a polynomial in $\zh$, then $\vartheta=\tau$, the rescaling given by the canonical coordinates of \cref{lem:set-canon} ($\tau\zh=-\frac54\zeta^iz/(u_i-u_i^{\mathrm c})$). This is not used, and not proved here.
\draftnote{Not proved here; the authors decide whether to keep this remark.}
That in the ansatz $M=\Mo(\vartheta\zh)$ the regularity of the one-edge graphs of genus $g$ determines $\vartheta$ modulo $u^{2g-2}$ is known only from computations, at genus $5$ (modulo $u^8$, \cref{sec:evi-graphs}).
\end{rmk}

\begin{rmk}[the vertex term]
\label{rmk:mai-vertexterm}
The vertex term is not the bare model potential. Without leaf insertions, the vertex graph of the flat sum contributes $d^{2g-2}\Fm_g(v)=c_g\hat v^{2-2g}$ with $\hat v=v/d$ (\cref{def:red-flatsum}, \cref{thm:mod-gap}). Here $d=(P_0)_{33}=1/p_{0,0}=\tau^{1/4}\varpi_0$, by symplecticity at $z^0$ ($(P_0)_{00}(P_0)_{33}=1$, \cref{lem:tra-T}~(d)) and \cref{lem:gau-diag}; so $\hat v=u\tau^{-3/2}/\varpi_0=u\tau^{-3/2}/\hat x$. The leaf insertions convert $\hat v$ into $\tc$: the full vertex term is $c_gT_{\Pc}^{2-2g}+\bbQ\fps{u}$ with $T_{\Pc}=\tc$ (\cref{thm:red-main}~(b), \cref{prop:tra-can}), and $\tc$ enters only as the vanishing coordinate $\Pi_0$ of the leaf point (\cref{lem:tra-T}). The first correction is discussed in \cref{rmk:red-leaf}.
\end{rmk}

\subsection{Universality}
\label{sec:mai-univ}

The flat half of the proof of Theorem~A, \cref{cor:red-can} with \cref{thm:mod-MLK}, uses the formal quintic only through its cyclic system and the canonical gauge; in flat form it applies to every self-dual cyclic system with one simple pole whose product is that of a pulled-back model. The comparison with the outer terms (\cref{thm:flh-main}) and the normalisation (\cref{thm:gam-D}) are specific to the formal quintic and are not part of this subsection. Throughout this subsection the systems are over $\bbQ$, in the sense of \cref{def:gau-onepole}, and for a non-zero series $y\in u\bbQ\fps{u}$ (a local letter), $\tilde A_y$ denotes the model pulled back by $y$, with coefficients $(y',-y'/y,y',-y',-y')$ in $\bbQ\lser{u}$ and product $-y'^5/y$ (\cref{prop:gau-equalprod}); it has the one-pole shape if and only if $y'(0)\neq0$.

\begin{lem}[the pole order]
\label{lem:mai-pole}
Let $A=\sum_ab_aE_{a,a+1}$ with $b_a\in\bbQ\lser{u}$ non-zero, $\Bprod\coloneqq\prod_ab_a$, and let $y\in u\bbQ\fps{u}$ be non-zero of order $n\geq1$.
\begin{enumerate}
\item If the gauge equation $z\,\partial_uP=AP-P\tilde A_y$ has a solution whose $z^0$-coefficient $P_0$ is non-zero, then $\Bprod=-y'^5/y$.
\item In that case $\ord_0\Bprod=4n-5\in\{-1,3,7,\dots\}$. If $\Bprod$ has a pole at $u=0$, then $n=1$, $\ord_0\Bprod=-1$ and $y'(0)^4=-\Res_{u=0}\Bprod$.
\end{enumerate}
\end{lem}

\begin{proof}
(i) At $z^0$ the gauge equation reads $AP_0=P_0\tilde A_y$, hence $A^5P_0=P_0\tilde A_y^5$; for a cyclic system $A^5=\Bprod\cdot1$, and $\tilde A_y^5=-(y'^5/y)\cdot1$. So $(\Bprod+y'^5/y)P_0=0$ over the field $\bbQ\lser{u}$, and $P_0\neq0$ gives the claim; this is the necessity clause of \cref{prop:gau-equalprod}. (ii) If $y$ has order $n$, then $y'^5$ has order $5(n-1)$ and $y'^5/y$ has order $5(n-1)-n=4n-5$. A pole forces $4n-5<0$, i.e. $n=1$; then $y=y'(0)u+O(u^2)$ and $y'^5/y=y'(0)^4u^{-1}+O(1)$, so $\Res_{u=0}\Bprod=-y'(0)^4$.
\end{proof}

The lemma is relative to the fixed model $\Wv$: a product with a pole of another order is not the product of a model pulled back by a power series.

\begin{lem}[a change of the product is a change of coordinate]
\label{lem:mai-coord}
Let $X$ be a self-dual one-pole system with coefficients $b_a$ and product $\Bprod_X$, and $v_X\in u\bbQ\fps{u}$ with $v_X'^5=-\Bprod_Xv_X$ and $a_X\coloneqq v_X'(0)\neq0$. Let $v_{\mathrm q}\coloneqq v(u)$ be the function of \cref{def:set-tau}, so that $-v_{\mathrm q}'^5/v_{\mathrm q}=-1/(u(1-u)^5)$ is the product of the quintic system (\cref{prop:gau-quintic}), and let $\phi\coloneqq v_X^{-1}\circ v_{\mathrm q}\in w\bbQ\fps{w}$, where $w$ is the new coordinate; $\phi'(0)=1/a_X$. Let $\tilde X$ be the system in the coordinate $w$ with coefficients $\tilde b_a(w)\coloneqq\phi'(w)\,b_a(\phi(w))$.
\begin{enumerate}
\item[(a)] $\tilde X$ is a self-dual one-pole system with the product of the quintic system, its model pull-back is $v_X\circ\phi=v_{\mathrm q}$, and it has the same residue as $X$.
\item[(b)] If $P(u,z)$ is a regular gauge from $\tilde A_{v_X}$ to $X$, then $P(\phi(w),z)$ is a regular gauge from $\tilde A_{v_{\mathrm q}}$ to $\tilde X$, with the same $P_{11}(0,z)$, $P_{00}(0,z)$ and $P_0(0)$; graded gauges go to graded gauges, and the canonical gauge to the canonical gauge.
\item[(c)] The vanishing solutions pull back: $Y^{\mathrm C}_{\tilde X}(w)=Y^{\mathrm C}_X(\phi(w))$, so (G) holds for $X$ if and only if it holds for $\tilde X$, and $T_{\tilde X}=T_X\circ\phi$, where $T_X\coloneqq\int_0^ub_0\,\dd u$.
\item[(d)] For a dressing $P$ of the model at $v_X(u)$ (\cref{def:tra-dressing}) and $\tilde P(w,z)\coloneqq P(\phi(w),z)$, a dressing of the model at $v_{\mathrm q}(w)$, every term of the flat sums satisfies $F^{\tilde X}_\Gamma(w)=F^X_\Gamma(\phi(w))$, where $F^X_\Gamma$ is the term of $\Gamma$ in the flat sum of $P.\Omega^{\mathrm{mod}}_{v_X(u)}$ and $F^{\tilde X}_\Gamma$ that of $\tilde P.\Omega^{\mathrm{mod}}_{v_{\mathrm q}(w)}$.
\item[(e)] A graph is regular for $X$ if and only if it is regular for $\tilde X$, and a vertex law $c_gT^{2-2g}+\bbQ\fps{u}$ transports to $c_g(T\circ\phi)^{2-2g}+\bbQ\fps{w}$.
\end{enumerate}
\end{lem}

\begin{proof}
(a) $\prod_a\tilde b_a=\phi'^5\Bprod_X(\phi)=-\phi'^5v_X'(\phi)^5/v_X(\phi)=-(v_X\circ\phi)'^5/(v_X\circ\phi)=-v_{\mathrm q}'^5/v_{\mathrm q}$. Each $\tilde b_a$ is $\phi'\,b_a\circ\phi$, so the self-dual shape $\tilde b_{a-1}=\tilde b_{3-a}$ is kept, and since $\phi(0)=0$ and $\phi'(0)\neq0$ the orders at $0$ are kept; the residue is $\lim_{w\to0}w\phi'(w)b_1(\phi(w))=\lim_{u\to0}ub_1(u)$, because $w\phi'(w)/\phi(w)\to1$. (b) Let $A_X=\sum_ab_aE_{a,a+1}$. By the chain rule, $z\,\partial_wP(\phi(w),z)=\phi'(w)\,(A_XP-P\tilde A_{v_X})(\phi(w))$, and the systems in the coordinate $w$ are $\phi'\cdot A_X\circ\phi$, which is $\tilde X$, and $\phi'\cdot\tilde A_{v_X}\circ\phi=\tilde A_{v_X\circ\phi}$. The coefficients stay in $\bbQ\fps{w}$, the grading is untouched, and the normalisations are conditions at the origin, which $\phi$ fixes. (c) The vanishing solutions are built by $y^{(n+1)}_a=\int_0^ub_a\,y^{(n)}_{a+1}\,\dd u$ from $y^{(0)}=-e_1$ (\cref{def:gau-van}), and $\int_0^{\phi(w)}h(u)\,\dd u=\int_0^wh(\phi(w'))\,\phi'(w')\,\dd w'$; the model's vanishing solution is a function of $v$ and $v_{\tilde X}=v_X\circ\phi$. (d) Every ingredient of the flat sum (the edges and the leaf point of $P$, the factor $d$, and the model correlators at $v_X(u)$, \cref{def:red-flatsum}) is formed pointwise in $u$, without derivatives in $u$, and every term is a finite sum (\cref{lem:tra-F}). (e) $f\mapsto f\circ\phi$ maps $\bbQ\lser{u}$ isomorphically onto $\bbQ\lser{w}$ and $\bbQ\fps{u}$ onto $\bbQ\fps{w}$.
\end{proof}

In a hatted normalisation, in which the term of a genus-$g$ graph is multiplied by $u^{2g-2}$ as in \cref{lem:flh-hat}, (d) reads $\tilde G_\Gamma(w)=(w/\phi(w))^{2g-2}G_\Gamma(\phi(w))$. By the lemma, a self-dual one-pole system whose product is that of a model pulled back by a power series, $-v_X'^5/v_X$ with $v_X\in u\bbQ\fps{u}$, is a coordinate change of a system with the product of the quintic system; the residue of the product is a coordinate effect as well (take $\phi$ linear).

\begin{cor}[universality, flat form]
\label{cor:mai-univ}
Let $v_A\in u\bbQ\fps{u}$ with $v_A'(0)\neq0$, and let $A$ be a self-dual one-pole system (\cref{def:gau-onepole}) with product $-v_A'^5/v_A$ and residue $r$; assume $r=-1/\lambda^2$ with $\lambda\in\bbQ^\times$ (a letter local to this corollary and its proof). Let $P^A_{\mathrm{can}}$ be the canonical gauge from $\tilde A_{v_A}$ to $A$, that is, the regular gauge with $(P^A_{\mathrm{can}})_{11}(0,z)=1$, which exists, is unique, is graded and has rational coefficients, as in \cref{def:gau-can} (\cref{thm:gau-U}, \cref{prop:gau-graded}); and let $P_A\coloneqq\lambda^{-1}P^A_{\mathrm{can}}$. Then for every $g\geq2$:
\begin{enumerate}
\item the term of every leg-free stable graph of genus $g$ with an edge in the flat sum of $P_A.\Omega^{\mathrm{mod}}_{v_A(u)}$ lies in $\bbQ\fps{u}$;
\item the term of the vertex graph lies in $c_g(T_A/\lambda)^{2-2g}+\bbQ\fps{u}$, where $T_A\coloneqq\int_0^ub_0\,\dd u$.
\end{enumerate}
\end{cor}

For the formal quintic in a conifold frame, $A$ is the system of \cref{prop:gau-quintic}, $v_A=v(u)$, $r=-1$ and $T_A=\tc$ (\cref{lem:gau-vanq}~(d)), and the corollary is \cref{cor:red-can} with \cref{thm:mod-MLK}. For a system $b=(b_0,\Bprod/(b_0^2b_3^2),b_0,b_3,b_3)$ with $\Bprod=-v_A'^5/v_A$ one has $r=-v_A'(0)^4/(b_0(0)b_3(0))^2$ by \cref{lem:mai-pole}, so $\lambda=\pm b_0(0)b_3(0)/v_A'(0)^2$; the vertex law is even in $\lambda$.

\begin{proof}
For the pair $(A,\tilde A_{v_A})$ the residue of $\tilde A_{v_A}$ is $\tilde r=-1$. By \cref{lem:gau-R1}, column $1$ of $P(0,z)$ is $P_{11}(0,z)e_1$ and $rP_{22}(0,z)=\tilde rP_{11}(0,z)$ for every regular gauge $P$. The proof of \cref{prop:gau-symp} uses the residues only in its last step, the evaluation of the entry $(2,1)$ at $u=0$ by \cref{lem:gau-R1}; with the general form of that lemma it gives $P(-z)^T\eta P(z)=(\tilde r/r)\,P_{11}(0,z)P_{11}(0,-z)\,\eta$. For $P_A$ one has $P_{11}(0,z)=\lambda^{-1}$ and $\tilde r/r=\lambda^2$, so $P_A$ is symplectic. By \cref{thm:gau-G}~(c), $P^A_{\mathrm{can}}\YC(v_A(u))=Y^{\mathrm C}_A$, since $\YC(v_A(u))$ is the vanishing solution of $\tilde A_{v_A}$ (\cref{lem:gau-vanmod}, the chain rule and \cref{lem:gau-uniq}, as in the proof of \cref{thm:gau-G}); so $P_A\YC(v_A(u))=\lambda^{-1}Y^{\mathrm C}_A$ has no positive power of $z$ (\cref{lem:gau-vanq}), which is (A), and $T_{P_A}=T_A/\lambda$ by \cref{lem:gau-vanq}~(d). $P_A$ is graded with coefficients in $\bbQ\fps{u}$ and invertible constant term (\cref{thm:gau-U}, \cref{prop:gau-graded}), so it is a dressing (\cref{def:tra-dressing}) with (Sp) and (A). \Cref{thm:red-main} with \cref{thm:mod-MLK} gives (i) and (ii).
\end{proof}

\begin{rmk}[non-gauge dressings]
\label{rmk:mai-annih}
\Cref{thm:red-main} does not use the gauge equation. It applies to every dressing $P$ with (Sp) and (A) in the sense of \cref{def:tra-dressing}, whether or not $P$ solves a gauge equation, and with \cref{thm:mod-MLK} it gives, for every $g\geq2$, that every edge graph of the flat sum of $P.\Omega^{\mathrm{mod}}_{v(u)}$ is regular and that the vertex graph lies in $c_gT_P^{2-2g}+\bbQ\fps{u}$. Here, as in \cref{cor:mai-univ}, only the flat form is claimed: a statement about a hatted sum would need an identity of the type of \cref{lem:flh-inputs}~(ii), which is available only for the rows of the formal quintic.
\end{rmk}

\begin{rmk}[a non-canonical normalisation]
\label{rmk:mai-stored}
Let $P^\sharp\coloneqq\Pc\,\cn(z)$, $\cn\in1+z^5\bbQ\fps{z^5}$, be the regular gauge of the other normalisation of \cref{sec:sgn-norm}, and write $\cn(z)=1+m_5z^5+m_{10}z^{10}+\cdots$ ($P^\sharp$ is a local letter). By \cref{thm:gau-U}, $P^\sharp$ is a graded regular gauge: it solves the gauge equation and is regular in $u$. By \cref{prop:gau-symp}, $P^\sharp(-z)^T\eta P^\sharp(z)=\cn(z)\cn(-z)\eta$ with $\cn(z)\cn(-z)=1-(m_5^2-2m_{10})z^{10}+\cdots$: $P^\sharp$ is symplectic modulo $z^{10}$, and symplecticity constrains only the even part of $\log\cn$. By \cref{thm:gau-G}, $P^\sharp\YC=\cn(z)\YCQ$, whose part of positive degree in $z$ begins with $m_5(z^5\YCQ)_{+}=-m_5(z^5e_1+\tc z^4e_0+\cdots)$; by \cref{thm:gau-G}~(b), $P^\sharp$ satisfies (A) only if $\cn=1$. So the gauge equation does not imply (A); this holds by construction, through the freedom $\cn(z^5)$ of \cref{thm:gau-U}, and it is (G) that fixes the odd part of $\log\cn$. The value of $m_5$, which is known only by computation, is recorded in \cref{sec:evi}, and the graphs that become singular in \cref{sec:evi-univ}.
\end{rmk}

\subsection{Remarks}
\label{sec:mai-rmk}

\begin{rmk}[graph-locality]
\label{rmk:mai-graphlocal}
The computations of \cref{sec:evi-univ} show that in the normalisation of \cref{rmk:mai-stored} six edge graphs of genus $4$ become singular, while the singular part of the total is unchanged: the normalisation only moves singular terms between the vertex graph and these edge graphs. So the singular part of the total does not determine the singular parts of the graph terms, and no argument that uses the total potential only through its singular part, such as the gap of the total, can prove Theorem~A. The proof is graph-local, and it uses the odd part of $\log\cn$, which is fixed by (G) (\cref{thm:gau-G}) and not by symplecticity.
\end{rmk}

\begin{rmk}[the hypotheses cannot be dropped]
\label{rmk:mai-necessity}
The computations of \cref{sec:evi-univ} give dressings that violate (A), among them $P^\sharp$ of \cref{rmk:mai-stored}, and dressings that are not regular in $u$ in the flat frame, whose flat sums have singular edge graphs: neither hypothesis of \cref{thm:red-main} can be dropped in general. For (Sp) no such example exists, because without it the numerator of the edge does not vanish at $\psi=-\psi'$ and the edge of \cref{def:tra-edge} is not a power series. The flatness in $u$, i.e. the gauge equation, is not a hypothesis (\cref{rmk:mai-annih}).
\end{rmk}

\begin{rmk}[the shape of the factorisation]
\label{rmk:mai-analogy}
\Cref{cor:mai-natfact} has the shape \emph{(a gauge regular at the singular point) $\times$ (the $R$-matrix of a universal local model in its natural variable)}. When the local model is a single critical point next to a logarithmic singularity, as for an integral of Gamma-function type with a constant weight, its $R$-matrix is a Stirling series in $z$ divided by that weight, and the second factor does not depend on $q$. Here the local model is $\Wv(p)=p^4/4+v/p$ with the form $\dd p$, whose five critical points coalesce onto the pole as $v\to0$; its natural variable $z/\varsigma^4=\tau(u)\zh$ is, up to the constant factor $-\frac54\zeta^i$, $z$ divided by the canonical coordinate measured from the conifold point (\cref{lem:set-canon}), which degenerates there, so the second factor depends on $u$. The counterpart of a scalar Hodge-type factor is $\Mo(\tau\zh)$ together with the regular gauge, not $\Mo$ alone. The form is natural, not forced: computations through genus $3$ show that the graph sums alone do not single it out.\draftnote{The computation behind the last sentence (the factorisations that make every edge graph of genus at most $3$ regular, besides the natural form) is not recorded in the evidence appendix; record it there, with its range, or drop the sentence.}
\end{rmk}


